\section{Task 4: Style-Conditioned Face-to-Sketch cGAN}\label{sec:t4}

\subsection{Model}
The generator $\hat{y}=G(x,s)$ (Fig.~\ref{fig:arch4}) is a U-Net~\cite{ronneberger2015}: seven stride-2 encoder blocks (64, 128, 256, 512 channels, four times 512 at the deepest levels) reduce the $128\times128$ photograph to a $1\times1$ code, and six transposed-convolution blocks with skip connections restore the resolution; a final transposed convolution with tanh produces a one-channel sketch. The first three decoder blocks use dropout (0.30). The style $s\in\{1,2,3\}$ is a \emph{learned categorical embedding} (three vectors of 32 values). It is incorporated into the generator twice, broadcast over the image as extra input channels and concatenated to the $1\times1$ bottleneck, and into the discriminator as extra input channels; it is therefore part of the model and not merely an interface label. The generator has 42.1\,M parameters. The discriminator is a PatchGAN~\cite{isola2017,li2016mgan}: it receives the photograph, the style map and either the real or the generated sketch (36 channels), and outputs a $14\times14$ map of patch logits (2.8\,M parameters).

\begin{figure*}[t]
\centering
\includegraphics[width=\textwidth]{arch_task4.pdf}
\caption{Style-conditioned cGAN: U-Net generator with a style embedding in the input and the bottleneck, and PatchGAN discriminator.}\label{fig:arch4}
\end{figure*}

\subsection{Losses and training}
The discriminator is trained with binary cross-entropy on logits to label real pairs $(x,y,s)$ as one and generated pairs $(x,G(x,s),s)$ as zero; its real and fake losses are recorded separately. The generator minimises
\begin{equation}
\mathcal{L}_G=\mathcal{L}_{adv}+\lambda_{L1}\,\mathcal{L}_{L1}\bigl(y,G(x,s)\bigr),
\end{equation}
the adversarial term being the discriminator's error on generated pairs. Both networks use Adam ($\beta_1=0.5$, $\beta_2=0.999$). Optuna tuned the generator and discriminator learning rates, batch size, base channel count, dropout, style-embedding size and $\lambda_{L1}$ with 15-epoch trials on all 899 training pairs (three completed, five pruned). The selected configuration (generator learning rate $1.5\times10^{-4}$, discriminator learning rate $8.6\times10^{-4}$, batch 8, 64 base channels, dropout 0.30, embedding size 32, $\lambda_{L1}=121$) was then trained for the full 100 epochs, recording the discriminator real and fake losses, the generator adversarial and L1 losses and the validation L1, SSIM and PSNR separately for every epoch (Fig.~\ref{fig:curves4}). Every ten epochs the same eight validation photographs are rendered in all three styles and logged to MLflow (Fig.~\ref{fig:prog}), so the development of the generator can be inspected over time.

\begin{figure*}[t]
\centering
\includegraphics[width=\textwidth]{curves_task4.pdf}
\caption{Task 4 training: discriminator and generator losses, and validation SSIM and L1 (fixed validation pairs).}\label{fig:curves4}
\end{figure*}

\begin{figure*}[t]
\centering
\includegraphics[width=\textwidth]{task4_progress.pdf}
\caption{The same eight validation photographs (top row of each panel) rendered by the generator in styles 1, 2 and 3 after epochs 1, 20 and 100.}\label{fig:prog}
\end{figure*}

\subsection{Results}
On the 1,046 official test pairs the selected model reaches L1 \nFourLOne{}, SSIM \nFourSSIM{} and PSNR \nFourPSNR{}\,dB (Table~\ref{tab:task4}, Fig.~\ref{fig:ex4}). The error differs strongly between styles: style~2 has an L1 error of \nFourLOnesTwo{} against \nFourLOnesOne{} for style~1 and \nFourLOnesThree{} for style~3, because its reference sketches use heavy dark strokes which a soft generator cannot match pixel by pixel, while style~1 and style~3 are lighter. The test styles are imbalanced (619, 381 and 46 pairs), so the style-3 figures rest on few samples. Each style is visibly different for the same photograph (Fig.~\ref{fig:styles}): style~2 is darkest and heaviest, style~1 lightest.

\begin{table}[t]
\centering
\caption{Task 4 on the official FS2K test set. The last row is a second training recipe (photo-only brightness/contrast jitter and learning-rate decay) that was tried and rejected because it was worse.}\label{tab:task4}
\small
\adjustbox{max width=\columnwidth}{\begin{tabular}{lrccc}
\toprule
Test subset & Pairs & L1 & SSIM & PSNR (dB) \\
\midrule
Style 1 & 619 & 0.0791 & 0.534$\pm$0.107 & 17.55 \\
Style 2 & 381 & 0.1522 & 0.395$\pm$0.082 & 12.47 \\
Style 3 & 46 & 0.0656 & 0.600$\pm$0.074 & 18.90 \\
\midrule
all (selected model) & 1046 & 0.1051 & 0.486 & 15.76 \\
all (rejected variant) & 1046 & 0.1112 & 0.451 & 15.23 \\
\bottomrule
\end{tabular}
}
\end{table}

\begin{figure*}[t]
\centering
\includegraphics[width=0.78\textwidth]{task4_examples.png}
\caption{Test examples: photograph, ground-truth sketch, generated sketch and absolute error (twelve pairs spread over the test set).}\label{fig:ex4}
\end{figure*}

\begin{figure}[t]
\centering
\includegraphics[width=\columnwidth]{task4_styles.png}
\caption{One photograph rendered in the three styles.}\label{fig:styles}
\end{figure}

\textbf{Why the sketches are soft.} Compared with the real sketches the generated ones are smoother (Fig.~\ref{fig:ex4}); we measured it. The average gradient magnitude of the generated sketches is only \nSharpShareAll\% of that of the real sketches (Table~\ref{tab:sharp}). A sharper input helps a little and a blurred input hurts a little (Table~\ref{tab:sens}), but the effect is small next to the 63\% gap, so the limit lies in the training rather than in the $128\times128$ input. Two causes are consistent with the evidence. With only 899 training pairs and $\lambda_{L1}=121$ the L1 term dominates and regresses towards an average stroke; and in the selected Optuna trial the discriminator stayed close to chance for the whole 15 epochs (real and fake loss near $\ln 2\approx0.69$), so the adversarial term contributed little; in the 100-epoch run it did learn to separate real from fake (losses fell to about 0.4 and the generator's adversarial loss rose to about 2, Fig.~\ref{fig:curves4}), but the generator's L1 term still dominates the objective. This is the perception--distortion trade-off~\cite{blau2018}: pixel metrics reward smooth outputs. A second recipe (photo-only brightness and contrast jitter, linear learning-rate decay) did not help (SSIM \nFourSSIMv{}, L1 \nFourLOnev{}) and was discarded. Stronger adversarial or perceptual losses, more data, or a higher resolution are the obvious next steps and were not explored.

\begin{table}[t]
\centering
\caption{Edge strength (mean gradient magnitude) of real and generated sketches on the FS2K test set.}\label{tab:sharp}
\small
\adjustbox{max width=\columnwidth}{\begin{tabular}{lrccc}
\toprule
Test subset & Pairs & Real sketch edge strength & Generated & Generated / real \\
\midrule
Style 1 & 619 & 0.0735 & 0.0267 & 36\% \\
Style 2 & 381 & 0.1506 & 0.0554 & 37\% \\
Style 3 & 46 & 0.0613 & 0.0368 & 60\% \\
\midrule
all & 1046 & 0.1011 & 0.0376 & 37\% \\
\bottomrule
\end{tabular}
}
\end{table}

\begin{table}[t]
\centering
\caption{Sensitivity of the generated edge strength to the sharpness of the input photograph.}\label{tab:sens}
\small
\adjustbox{max width=\columnwidth}{\begin{tabular}{lcc}
\toprule
Input photo (300 test faces) & Generated edge strength & L1 to real sketch \\
\midrule
input blurred (Gaussian, r=1.5) & 0.0346 & 0.1108 \\
input as is & 0.0381 & 0.1084 \\
input sharpened (unsharp mask) & 0.0436 & 0.1091 \\
\bottomrule
\end{tabular}
}
\end{table}

\begin{figure}[t]
\centering
\includegraphics[width=\columnwidth]{task4_failures.png}
\caption{The four test pairs with the lowest SSIM.}\label{fig:fail4}
\end{figure}

\textbf{Failure cases.} Figure~\ref{fig:fail4} shows the four test pairs with the lowest SSIM (0.18--0.21). Three are style~2 portraits with long hair: the reference sketches draw the hair with dense, dark, hand-drawn strokes, while the generator renders it as a smooth grey mass and draws the face softly (SSIM 0.18, 0.20, 0.20). The fourth is a style~1 portrait with long hair and bangs, where the reference consists of thin, crisp lines and the generated sketch is pale and blurred (0.21). The faces remain recognisable, so these are failures of stroke detail, not of identity. The error maps light up along every reference stroke: a thin line that is missing or displaced by a few pixels costs a large pixel error, which is why pixel metrics punish a soft generator so heavily on hair-rich, style~2 subjects.
